\documentclass[11pt,a4paper]{article}

\usepackage[utf8]{inputenc}
\usepackage[T1]{fontenc}
\usepackage{lmodern}
\usepackage{amsmath,amssymb,amsthm}
\usepackage{geometry}
\geometry{margin=2.5cm}
\usepackage{array}

% Rule commands (booktabs not available)
\newcommand{\toprule}{\hline\hline}
\newcommand{\midrule}{\hline}
\newcommand{\bottomrule}{\hline\hline}

\setlength{\parindent}{0pt}
\setlength{\parskip}{0.5em}

\newtheorem{definition}{Definition}
\newtheorem{proposition}{Proposition}
\newtheorem{remark}{Remark}

\title{\textbf{Information Coefficient (IC) Reference}\\[0.5em]
\large Mathematical Foundations for Signal Quality Assessment\\[0.3em]
\normalsize NAT Quantitative Research Platform}
\author{NAT Research}
\date{June 2026}

\begin{document}
\maketitle
\tableofcontents
\newpage

%===========================================================================
\section{Introduction}
%===========================================================================

The Information Coefficient (IC) is the canonical metric for measuring a signal's
predictive power over future asset returns. In the NAT platform, IC governs every
stage of the research pipeline: discovery, hypothesis gating, walk-forward
validation, regime conditioning, and live decay monitoring.

This document provides the closed-form definitions, derived metrics, statistical
tests, and threshold configurations as implemented in the NAT codebase. It is
intended as a self-contained reference for developing new IC-based algorithms
targeting Hyperliquid perpetual futures (BTC, ETH, SOL) across tick-level to
macro horizons.

\medskip
\textbf{Source files:}
\begin{itemize}    \item \texttt{scripts/exploration/spannung\_regime\_screener.py} --- core IC computation
    \item \texttt{scripts/alpha/screener.py} --- rolling/expanding IC, FDR correction
    \item \texttt{scripts/agent/gates.py} --- gate protocol (G1, G2)
    \item \texttt{scripts/agent/base.py} --- decay monitoring
    \item \texttt{scripts/alpha/regime\_filter.py} --- conditional IC
    \item \texttt{config/agent.toml} --- threshold configuration
\end{itemize}

%===========================================================================
\section{Core Definitions}
%===========================================================================

%-----------------------------------
\subsection{Forward Log Return}
%-----------------------------------

\begin{definition}[Forward Return]
Given a price series $\{p_t\}$ and a horizon $h$ (in ticks or bars), the
forward log return at time $t$ is:
\begin{equation}
    r_{t \to t+h} = \ln\!\left(\frac{p_{t+h}}{\max(p_t,\; \varepsilon)}\right)
    \label{eq:fwd_return}
\end{equation}
where $\varepsilon = 10^{-15}$ prevents division by zero. The final $h$
entries are set to NaN (no forward return beyond the array boundary).
\end{definition}

\begin{remark}
The horizon $h$ is expressed in native sampling units. At tick level
(100\,ms emission), $h = 10$ corresponds to a 1\,s horizon; $h = 50$
corresponds to 5\,s. For bar-resampled data, $h$ is in bars.
\end{remark}

%-----------------------------------
\subsection{Spearman Rank IC}
%-----------------------------------

\begin{definition}[Information Coefficient]
The IC of a signal $S$ at horizon $h$ is the Spearman rank correlation
between the signal value at time $t$ and the forward return:
\begin{equation}
    \mathrm{IC} = \rho_s\!\left(S_t,\; r_{t \to t+h}\right)
    = 1 - \frac{6 \sum_{i=1}^{n} d_i^2}{n(n^2 - 1)}
    \label{eq:spearman}
\end{equation}
where $d_i = \mathrm{rank}(S_i) - \mathrm{rank}(r_i)$ and $n$ is the number
of valid (non-NaN) paired observations.
\end{definition}

\begin{remark}
Spearman rank correlation is preferred over Pearson because:
(i)~it is robust to outliers common in microstructure data,
(ii)~it captures monotonic (not just linear) relationships,
(iii)~it is invariant to monotone transformations of the signal.
The sole exception in NAT is the Kalman filter tuning module, which uses
Pearson correlation for filter parameter optimization.
\end{remark}

%-----------------------------------
\subsection{Rolling Non-Overlapping IC}
%-----------------------------------

In practice, IC is computed over rolling non-overlapping windows to ensure
statistical independence between estimates.

Given $N$ observations, partition into $K = \lfloor N / W \rfloor$ windows
of size $W$. For window $k$:
\begin{equation}
    \mathrm{IC}_k = \rho_s\!\left(
        S_{kW:(k+1)W},\;
        r_{kW:(k+1)W}
    \right)
    \label{eq:rolling_ic}
\end{equation}

A window is \textbf{valid} if and only if:
\begin{enumerate}    \item Number of finite paired observations $\geq n_{\min}$ (default: 100)
    \item Signal range $\mathrm{ptp}(S) > 10^{-15}$ (not constant)
    \item Return range $\mathrm{ptp}(r) > 10^{-15}$ (not constant)
    \item Spearman computation does not return NaN
\end{enumerate}

\textbf{Default parameters:} $W = 3000$ ticks, $n_{\min} = 100$.

%===========================================================================
\section{Derived Metrics}
%===========================================================================

Given $K$ valid IC windows $\{\mathrm{IC}_1, \ldots, \mathrm{IC}_K\}$:

%-----------------------------------
\subsection{IC Mean and Standard Deviation}
%-----------------------------------

\begin{align}
    \overline{\mathrm{IC}} &= \frac{1}{K} \sum_{k=1}^{K} \mathrm{IC}_k
    \label{eq:ic_mean} \\[6pt]
    \sigma_{\mathrm{IC}} &= \sqrt{\frac{1}{K} \sum_{k=1}^{K}
        \left(\mathrm{IC}_k - \overline{\mathrm{IC}}\right)^2}
    \label{eq:ic_std}
\end{align}

%-----------------------------------
\subsection{Information Ratio (IC IR)}
%-----------------------------------

\begin{definition}[Information Ratio]
The IC IR measures signal-to-noise of predictability:
\begin{equation}
    \mathrm{IR} = \frac{\overline{\mathrm{IC}}}{\sigma_{\mathrm{IC}}}
    \label{eq:ic_ir}
\end{equation}
\end{definition}

A signal with $\mathrm{IR} > 1$ has mean IC exceeding its variability ---
a strong indicator of persistent predictive power.

%-----------------------------------
\subsection{Annualized IC IR}
%-----------------------------------

To compare signals across timeframes, the IR is annualized:
\begin{equation}
    \mathrm{IR}_{\mathrm{ann}} = \mathrm{IR} \times
        \sqrt{\frac{B_{\mathrm{year}}}{B_{\mathrm{step}}}}
    \label{eq:ann_ir}
\end{equation}
where $B_{\mathrm{year}}$ is the number of bars per year and
$B_{\mathrm{step}} = N_{\mathrm{bars}} / K$ is the average bars per
IC window.

\begin{center}
\begin{tabular}{lrl}
    \toprule
    \textbf{Timeframe} & $B_{\mathrm{year}}$ & \textbf{Computation} \\
    \midrule
    15\,min & 35{,}040 & $96 \times 365$ \\
    1\,h    &  8{,}760 & $24 \times 365$ \\
    4\,h    &  2{,}190 & $6 \times 365$ \\
    \bottomrule
\end{tabular}
\end{center}

%-----------------------------------
\subsection{Statistical Significance}
%-----------------------------------

\subsubsection{t-Statistic and p-Value (Window-Based)}

For the null hypothesis $H_0\!: \overline{\mathrm{IC}} = 0$:
\begin{align}
    t &= \frac{\overline{\mathrm{IC}}}{\sigma_{\mathrm{IC}} / \sqrt{K}}
    \label{eq:t_stat} \\[6pt]
    p &= 2\,\Pr\!\left(|T| > |t|\right), \quad T \sim t_{K-1}
    \label{eq:p_value}
\end{align}
where $t_{K-1}$ is the Student's $t$-distribution with $K-1$ degrees of freedom.

\subsubsection{Z-Test p-Value (Observation-Based)}

For a single IC estimate over $n$ observations, under $H_0\!: \mathrm{IC} = 0$:
\begin{equation}
    p = \mathrm{erfc}\!\left(\frac{|\mathrm{IC}| \cdot \sqrt{n}}{\sqrt{2}}\right)
    \label{eq:z_pvalue}
\end{equation}
where $\mathrm{erfc}$ is the complementary error function. This approximation
treats the Fisher-transformed Spearman correlation as asymptotically normal.

%-----------------------------------
\subsection{IC Hit Rate}
%-----------------------------------

\begin{equation}
    \mathrm{HitRate} = \frac{1}{K} \sum_{k=1}^{K}
        \mathbf{1}\!\left[\mathrm{sign}(\mathrm{IC}_k)
        = \mathrm{sign}(\overline{\mathrm{IC}})\right]
    \label{eq:hit_rate}
\end{equation}

A hit rate above 70\% indicates the signal consistently predicts in the same
direction across windows.

%-----------------------------------
\subsection{IC Autocorrelation}
%-----------------------------------

Lag-1 autocorrelation of the IC series:
\begin{equation}
    \rho_1 = \frac{\sum_{k=1}^{K-1}
        (\mathrm{IC}_k - \overline{\mathrm{IC}})
        (\mathrm{IC}_{k+1} - \overline{\mathrm{IC}})}
        {\sum_{k=1}^{K} (\mathrm{IC}_k - \overline{\mathrm{IC}})^2}
    \label{eq:ic_autocorr}
\end{equation}

Low autocorrelation ($|\rho_1| < 0.3$) suggests IC is noisy window-to-window;
high autocorrelation suggests persistent regimes of predictability.

%-----------------------------------
\subsection{Trading Cost Metrics}
%-----------------------------------

\subsubsection{Signal Turnover}

\begin{equation}
    \mathrm{Turnover} = \frac{\mathrm{mean}(|\Delta S_t|)}{\sigma(S)}
    \label{eq:turnover}
\end{equation}
where $\Delta S_t = S_t - S_{t-1}$. High turnover implies frequent
rebalancing and higher transaction costs.

\subsubsection{Breakeven Fee (bps)}

The minimum fee level at which the signal remains profitable:
\begin{equation}
    \mathrm{Breakeven}_{\mathrm{bps}} =
        \frac{|\overline{\mathrm{IC}}| \times \sigma(r)}
             {\mathrm{Turnover}} \times 10{,}000
    \label{eq:breakeven}
\end{equation}

\begin{remark}
On Hyperliquid, taker fees are $\sim$3.5\,bps and maker fees are
$\sim$0.2\,bps. A signal is cost-viable if $\mathrm{Breakeven}_{\mathrm{bps}}$
exceeds the relevant fee tier. When $\mathrm{Turnover} \approx 0$, the
breakeven is $+\infty$ (the signal never trades).
\end{remark}

%===========================================================================
\section{Computation Variants}
%===========================================================================

%-----------------------------------
\subsection{Rolling IC (Non-Overlapping)}
%-----------------------------------

The standard method (Section~2.3). Windows are non-overlapping to ensure
independence of IC estimates. Used in the spannung regime screener and
alpha screener.

\textbf{Parameters:} window size $W$, minimum observations $n_{\min}$.

%-----------------------------------
\subsection{Expanding IC (Causal)}
%-----------------------------------

To avoid lookahead bias in walk-forward evaluation, IC is computed on an
expanding window anchored at $t = 0$:

\begin{equation}
    \mathrm{IC}^{\mathrm{exp}}_j = \rho_s\!\left(
        S_{0:t_j},\; r_{0:t_j}
    \right), \quad
    t_j = n_{\min} + j \cdot \Delta t
    \label{eq:expanding_ic}
\end{equation}

where $\Delta t = \max(\lfloor N/20 \rfloor, 1)$ is the step size and
$n_{\min} = 50$ is the minimum data before the first IC estimate.

This produces a monotonically growing sample size, simulating what a
researcher would observe in real time.

%-----------------------------------
\subsection{Conditional IC (Regime-Gated)}
%-----------------------------------

For a binary regime mask $M_t \in \{0, 1\}$, the conditional IC is:

\begin{equation}
    \mathrm{IC}_{\mathrm{cond}} = \rho_s\!\left(
        S_t \cdot M_t,\; r_{t \to t+h} \cdot M_t
    \right)
    \label{eq:conditional_ic}
\end{equation}

computed only on observations where $M_t = 1$. The window size adapts:
\begin{equation}
    W_{\mathrm{cond}} = \min\!\left(W,\;
        \max\!\left(\left\lfloor\frac{n_{\mathrm{valid}}}{4}\right\rfloor,
        \; n_{\min}\right)\right)
    \label{eq:adaptive_window}
\end{equation}

If fewer than 2 valid windows exist, a single global Spearman IC is
computed as a fallback.

\textbf{Minimum regime observations:} $n_{\mathrm{regime}} \geq 500$.

%-----------------------------------
\subsection{Bandpass-Filtered IC}
%-----------------------------------

Signals are optionally filtered through an ultra-low bandpass before IC
computation:
\begin{equation}
    \tilde{S}_t = \mathrm{BPF}(S_t;\; f_{\mathrm{lo}}, f_{\mathrm{hi}})
    \label{eq:bandpass}
\end{equation}

\textbf{Default band:} $f_{\mathrm{lo}} = 0.005$\,Hz,
$f_{\mathrm{hi}} = 0.1$\,Hz. At 10\,Hz sampling (100\,ms ticks), this
extracts the slow-moving component with periods between 10\,s and 200\,s.

Filtered IC often exceeds raw IC by removing high-frequency noise while
preserving the predictive component.

%-----------------------------------
\subsection{Tick-Level vs.\ Bar-Resampled Paths}
%-----------------------------------

IC computation supports two data representations:

\begin{center}
\begin{tabular}{lll}
    \toprule
    \textbf{Path} & \textbf{Forward Return} & \textbf{Horizon} \\
    \midrule
    Tick-level & $r_t = (p_{t+h} - p_t) / p_t$ &
        $h = \max(1, \lfloor h_s / 0.1 \rfloor)$ ticks \\
    Bar-resampled & $r_t = (p_{t+h} - p_t) / p_t$ &
        $h = \max(1, \lfloor h_s / T_{\mathrm{bar}} \rfloor)$ bars \\
    \bottomrule
\end{tabular}
\end{center}

where $h_s$ is the horizon in seconds and $T_{\mathrm{bar}}$ is the bar
period in seconds.

%===========================================================================
\section{Statistical Quality Control}
%===========================================================================

%-----------------------------------
\subsection{Benjamini--Hochberg FDR Correction}
%-----------------------------------

When testing $m$ hypotheses simultaneously, false discovery rate (FDR)
control prevents spurious signals from surviving.

\textbf{Algorithm (Benjamini--Hochberg, 1995):}
\begin{enumerate}    \item Sort raw p-values: $p_{(1)} \leq p_{(2)} \leq \cdots \leq p_{(m)}$
    \item Compute adjusted p-values:
        \begin{equation}
            p_{(i)}^{\mathrm{adj}} = p_{(i)} \cdot \frac{m}{i}
            \label{eq:bh_adjust}
        \end{equation}
    \item Enforce monotonicity (step-down):
        \begin{equation}
            p_{(i)}^{\mathrm{adj}} \leftarrow
                \min\!\left(p_{(i)}^{\mathrm{adj}},\;
                p_{(i+1)}^{\mathrm{adj}}\right),
            \quad i = m-1, \ldots, 1
            \label{eq:bh_monotone}
        \end{equation}
    \item Clip to $[0, 1]$
    \item Reject $H_i$ where $p_i^{\mathrm{adj}} < q$
\end{enumerate}

\textbf{Default FDR threshold:} $q = 0.05$.

\begin{remark}
In the agent protocol, FDR is applied at the end of each research cycle
across all hypotheses that passed individual gates. This controls the
expected fraction of false discoveries in the promoted signal registry.
\end{remark}

%-----------------------------------
\subsection{5-Gate Hypothesis Protocol}
%-----------------------------------

Every candidate signal must pass a sequence of quality gates before
promotion. The IC-related gates are:

\medskip

\textbf{Gate 1 --- IC Gate:}
\begin{equation}
    \text{PASS} \iff |\mathrm{IC}| \geq \mathrm{IC}_{\min}
    \label{eq:gate1}
\end{equation}

The p-value is computed via~\eqref{eq:z_pvalue} for informational
reporting but does not gate independently.

\medskip

\textbf{Gate 2 --- $\Delta$IC Gate:}
\begin{equation}
    \Delta\mathrm{IC} = \mathrm{IC}_{\mathrm{gated}}
        - \mathrm{IC}_{\mathrm{baseline}}, \quad
    \text{PASS} \iff \Delta\mathrm{IC} \geq \Delta\mathrm{IC}_{\min}
    \label{eq:gate2}
\end{equation}

This ensures regime conditioning adds genuine value over the ungated
signal. Skipped when no regime gate is specified.

\medskip

\textbf{Gate 5 --- Regime Improvement:}
\begin{equation}
    \text{PASS} \iff \exists\; r \text{ s.t. }
        \frac{|\mathrm{IC}_r|}{|\mathrm{IC}_{\mathrm{global}}|} > 1.5
    \label{eq:gate5}
\end{equation}

At least one regime must improve IC by 50\% over the global (all-data)
baseline.

\medskip

\textbf{Full gate sequence:}
\begin{center}
\begin{tabular}{clll}
    \toprule
    \textbf{Gate} & \textbf{Name} & \textbf{Metric} & \textbf{Threshold} \\
    \midrule
    G1 & IC & $|\mathrm{IC}|$ & $\geq \mathrm{IC}_{\min}$ \\
    G2 & $\Delta$IC & $\mathrm{IC}_{\mathrm{gated}} - \mathrm{IC}_{\mathrm{base}}$
        & $\geq \Delta\mathrm{IC}_{\min}$ \\
    G3 & Cost & avg return per trade (bps) & $\geq 0.1$\,bps \\
    G4 & Coverage & regime observation fraction & $\geq 0.10$ \\
    G5 & Regime & $\max(|\mathrm{IC}_r| / |\mathrm{IC}_g|)$ & $> 1.5$ \\
    \bottomrule
\end{tabular}
\end{center}

After all gates, surviving hypotheses undergo BH-FDR correction
(Section~5.1) at $q = 0.05$.

%-----------------------------------
\subsection{Adaptive IC Threshold}
%-----------------------------------

As the signal registry matures, the acceptance bar tightens:
\begin{equation}
    \mathrm{IC}_{\min}^{\mathrm{adapt}} = \max\!\left(
        \mathrm{IC}_{\mathrm{floor}},\;
        0.8 \times \mathrm{median}\!\left(
            \{\mathrm{IC}_{\mathrm{exp}}^{(i)}\}_{i \in \mathrm{registry}}
        \right)
    \right)
    \label{eq:adaptive_ic}
\end{equation}

This prevents the registry from accumulating marginal signals once
strong ones have been found.

%===========================================================================
\section{IC Decay Monitoring}
%===========================================================================

Live signals are monitored for IC degradation. The decay protocol:

\begin{enumerate}    \item Compute rolling IC on the latest 2 days of data
    \item Compare against expected IC at registration:
        \begin{equation}
            \mathrm{IC}_{\mathrm{threshold}} =
                \mathrm{IC}_{\mathrm{expected}} \times \delta
            \label{eq:decay_threshold}
        \end{equation}
        where $\delta = 0.5$ (decay ratio).
    \item Track consecutive days below threshold:
        \begin{equation}
            D_t = \begin{cases}
                D_{t-1} + 1 & \text{if } \mathrm{IC}_{\mathrm{rolling}} <
                    \mathrm{IC}_{\mathrm{threshold}} \\
                0 & \text{otherwise (recovery)}
            \end{cases}
            \label{eq:decay_counter}
        \end{equation}
    \item Retire signal when $D_t \geq D_{\max}$ (default: 14 days)
\end{enumerate}

\textbf{Rolling IC computation:}
\begin{itemize}    \item Tick-level path: horizon in 100\,ms ticks,
        $h = \max(1, \lfloor h_s / 0.1 \rfloor)$
    \item Bar-resampled path: aggregate to bars first (5\,min or 1\,h),
        then compute Spearman IC
    \item Minimum valid data points: 500
    \item IC history: rolling buffer of 30 daily entries
\end{itemize}

%===========================================================================
\section{Regime-Conditioned Feature Weights}
%===========================================================================

When a regime improves IC, feature weights are adjusted per-regime:

\begin{equation}
    w_f^{(r)} = \begin{cases}
        \displaystyle\frac{|\mathrm{IC}_f^{(r)}|}
            {\sum_{f'} |\mathrm{IC}_{f'}^{(r)}|}
        & \text{if } \displaystyle\frac{|\mathrm{IC}_f^{(r)}|}
            {|\mathrm{IC}_f^{(\mathrm{global})}|} > 1.5 \\[12pt]
        \displaystyle\frac{|\mathrm{IC}_f^{(\mathrm{global})}|}
            {\sum_{f'} |\mathrm{IC}_{f'}^{(\mathrm{global})}|}
        & \text{otherwise}
    \end{cases}
    \label{eq:regime_weights}
\end{equation}

This allows the system to exploit regime-specific predictability while
falling back to global weights when conditioning does not help.

%===========================================================================
\section{Threshold Reference}
%===========================================================================

\subsection{IC Computation Parameters}

\begin{center}
\begin{tabular}{llr}
    \toprule
    \textbf{Parameter} & \textbf{Context} & \textbf{Value} \\
    \midrule
    \texttt{IC\_WINDOW}     & Rolling window size (ticks) & 3{,}000 \\
    \texttt{IC\_MIN\_OBS}   & Min valid obs per window    & 100 \\
    \texttt{MIN\_REGIME\_OBS} & Min obs per regime condition & 500 \\
    \texttt{ULTRALOW\_BAND} & Bandpass filter (Hz)        & $(0.005,\; 0.1)$ \\
    \texttt{FWD\_HORIZONS}  & Forward return horizons     & 1\,s = 10t,\; 5\,s = 50t \\
    Expanding min\_obs      & First expanding IC point    & 50 \\
    Expanding step          & Step size                   & $\max(\lfloor N/20 \rfloor, 1)$ \\
    Rolling min\_obs (bars) & Alpha screener              & 30 \\
    \bottomrule
\end{tabular}
\end{center}

\subsection{Agent Gate Thresholds}

\begin{center}
\begin{tabular}{lccc}
    \toprule
    \textbf{Parameter} & \textbf{Micro} & \textbf{MF (5\,min)} & \textbf{Macro (1\,h)} \\
    \midrule
    $\mathrm{IC}_{\min}$          & 0.10 & 0.08 & 0.07 \\
    $\Delta\mathrm{IC}_{\min}$    & 0.05 & 0.03 & 0.02 \\
    Min coverage                   & 0.10 & ---  & ---  \\
    Walk-forward sign consistency  & 0.80 & ---  & ---  \\
    \bottomrule
\end{tabular}
\end{center}

\textbf{Rationale:} Lower thresholds at longer horizons compensate for
fewer observations per day ($\sim$24 bars/day at 1\,h vs.\ $\sim$36{,}000
ticks/hour at tick level).

\subsection{Decay and FDR Parameters}

\begin{center}
\begin{tabular}{llr}
    \toprule
    \textbf{Parameter} & \textbf{Description} & \textbf{Value} \\
    \midrule
    $\delta$ (decay ratio) & Retire if rolling IC $< \delta \times$ expected & 0.5 \\
    $D_{\max}$ & Consecutive days before retirement & 14 \\
    $q$ (FDR) & Benjamini--Hochberg threshold & 0.05 \\
    OOS min dates & Minimum out-of-sample periods & 2 \\
    Min symbols & Cross-asset replication requirement & 2 \\
    \bottomrule
\end{tabular}
\end{center}

%===========================================================================
\section{Algorithm Development Guide}
%===========================================================================

%-----------------------------------
\subsection{Computing IC for a New Signal}
%-----------------------------------

To evaluate a candidate signal $S$:

\begin{enumerate}
    \item \textbf{Compute forward returns} at target horizon using~\eqref{eq:fwd_return}.
    \item \textbf{Rolling IC}: partition into non-overlapping windows of
        size $W$ and compute Spearman IC per window via~\eqref{eq:rolling_ic}.
    \item \textbf{Aggregate}: compute $\overline{\mathrm{IC}}$,
        $\sigma_{\mathrm{IC}}$, IR, t-stat, p-value.
    \item \textbf{Cost check}: compute turnover~\eqref{eq:turnover} and
        breakeven~\eqref{eq:breakeven}. Reject if breakeven $<$ 3.5\,bps
        (taker) or $<$ 0.2\,bps (maker).
    \item \textbf{FDR}: if testing multiple signals, apply BH correction
        at $q = 0.05$ before declaring significance.
\end{enumerate}

%-----------------------------------
\subsection{Gate Checklist}
%-----------------------------------

Before promoting a signal to the registry, verify:

\begin{enumerate}    \item[$\square$] G1: $|\mathrm{IC}| \geq \mathrm{IC}_{\min}$ for the
        target agent tier
    \item[$\square$] G2: $\Delta\mathrm{IC} \geq \Delta\mathrm{IC}_{\min}$
        (if regime-gated)
    \item[$\square$] G3: Average return per trade $\geq$ 0.1\,bps after costs
    \item[$\square$] G4: Regime coverage $\geq$ 10\% of observations
    \item[$\square$] G5: At least one regime with IC ratio $> 1.5\times$ global
    \item[$\square$] FDR: Survives BH correction across the hypothesis batch
    \item[$\square$] Cross-symbol: IC replicates on $\geq$ 2 symbols
    \item[$\square$] Walk-forward: $\geq$ 80\% sign consistency across OOS periods
\end{enumerate}

%-----------------------------------
\subsection{Interpreting IC Values}
%-----------------------------------

\begin{center}
\begin{tabular}{cl}
    \toprule
    $|\mathrm{IC}|$ & \textbf{Interpretation} \\
    \midrule
    $< 0.02$ & Noise --- no signal \\
    $0.02$--$0.05$ & Weak --- may survive in ensemble only \\
    $0.05$--$0.10$ & Moderate --- viable with low costs \\
    $0.10$--$0.20$ & Strong --- standalone signal candidate \\
    $> 0.20$ & Exceptional --- verify not overfit \\
    \bottomrule
\end{tabular}
\end{center}

\begin{remark}
IC values above 0.20 at tick level should be treated with suspicion.
Common causes: lookahead bias, stale feature lagging price, or
insufficient data. Always verify with expanding IC and cross-symbol
replication before trusting high IC values.
\end{remark}

\end{document}
